\section{Coherencia metodológica del proyecto}
\label{sec:coherencia}

Esta sección analiza la coherencia interna del estudio, verificando que el
diseño metodológico efectivamente responde a la pregunta de investigación
formulada.

\subsection{Cadena de articulación}

El proyecto exhibe una articulación vertical unívoca. El \textbf{fenómeno}
—degradación estocástica de la telemetría en chancado primario— se delimita como
\textbf{problema} informático de observabilidad parcial; de ese problema se
deriva la \textbf{pregunta} sobre la capacidad de un estimador estocástico para
reconstruir el estado interno; la pregunta se descompone en cuatro
\textbf{objetivos} que cubren caracterización, modelación, implementación y
validación; y cada objetivo determina una \textbf{decisión metodológica} con su
instrumento correspondiente. La tabla~\ref{tab:coherencia} formaliza esa cadena
e incorpora la evidencia efectivamente producida.

\begin{table}[H]
\centering
\caption{Matriz de coherencia metodológica y trazabilidad del proyecto.}
\label{tab:coherencia}
\footnotesize
\begin{tabular}{@{}p{2.9cm}p{3.2cm}p{3.2cm}p{3.9cm}@{}}
\toprule
\textbf{Objetivo específico} & \textbf{Decisión metodológica} &
\textbf{Instrumento} & \textbf{Evidencia producida}\\
\midrule
\textbf{OE1.} Caracterizar el ruido SCADA en chancado primario &
Modelación de ruido gaussiano mediante $\mathbf{Q}$ y $\mathbf{R}$
parametrizadas &
Revisión de literatura indexada y generador de degradación &
Tabla~\ref{tab:parametros}; verificación de que
$\mathrm{MSE}_{\text{SCADA}} \approx R$\\
\addlinespace
\textbf{OE2.} Modelar la dinámica de flujo en la tolva &
Espacio de estados lineal en tiempo discreto con estado latente &
Ecuaciones de balance de masa y matrices
$\mathbf{A},\mathbf{B},\mathbf{H}$ &
Ecuaciones~\eqref{eq:vector-estado}--\eqref{eq:matrices};
observabilidad verificada, ec.~\eqref{eq:observabilidad}\\
\addlinespace
\textbf{OE3.} Desarrollar el Filtro de Kalman Lineal &
Programación modular científica con núcleo desacoplado &
Ciclo recursivo predictor-corrector en forma de Joseph &
Módulo ejecutable; figura~\ref{fig:convergencia};
sección~\ref{sec:extensibilidad}\\
\addlinespace
\textbf{OE4.} Validar el desempeño frente a la lectura degradada &
Monte Carlo e inferencia estadística pareada &
Métricas MSE/RMSE, prueba t, Wilcoxon y análisis de sensibilidad &
Tablas~\ref{tab:resultados-volumen}--\ref{tab:degradacion};
figuras~\ref{fig:montecarlo}--\ref{fig:sensibilidad}\\
\bottomrule
\end{tabular}
\end{table}

\subsection{Verificación de que el diseño responde a la pregunta}

La pregunta de investigación exige tres cosas y el diseño provee evidencia para
cada una.

Exige, primero, \textbf{mitigar el ruido instrumental}. El diseño lo verifica
comparando el error de la señal filtrada contra la línea base declarada, con una
verdad de referencia conocida por construcción, condición que solo la simulación
permite garantizar.

Exige, segundo, \textbf{reconstruir el estado dinámico interno} en un sistema
\textbf{parcialmente observable}. El diseño aborda esta exigencia de manera
directa mediante la incorporación de la tasa de llenado como estado latente no
sensado, y verifica formalmente la condición de observabilidad que hace posible
su reconstrucción. Sin ese segundo estado, el estudio no podría afirmar nada
sobre observabilidad parcial: se limitaría a suavizar una señal medida.

Exige, tercero, que la reconstrucción sea \textbf{confiable}. El diseño
operacionaliza esa exigencia mediante una prueba de hipótesis formal con nivel
de significancia declarado, un intervalo de confianza y un contraste de robustez
no paramétrico, de modo que la afirmación de confiabilidad sea verificable y no
argumentativa.

\subsection{Ajustes realizados para fortalecer la coherencia}

La coherencia del diseño no fue estática: se fortaleció mediante ajustes
concretos durante el desarrollo, la mayoría de ellos derivados de la
retroalimentación docente y documentados en el Anexo~\ref{anexo:ajustes}.

Tres de esos ajustes merecen mención en este análisis por su impacto en la
coherencia interna. El primero es la \textbf{separación dimensional} entre MSE y
RMSE, que resolvió la inconsistencia de reportar métricas en
\si{\metre\tothe{6}} junto a intervalos de confianza en \si{\metre\cubed}, y que
condujo a la decisión de fundar toda la inferencia sobre el RMSE. El segundo es
la \textbf{verificación empírica del supuesto de normalidad}, que reveló su
incumplimiento y motivó la incorporación del contraste de Wilcoxon; sin ese
ajuste, la conclusión estadística habría descansado sobre un supuesto no
satisfecho. El tercero es la \textbf{recalibración de la covarianza de proceso}
descrita en la sección~\ref{sec:recalibracion}, que restituyó la
correspondencia entre el modelo y la física del activo.

Corresponde señalar que los dos primeros ajustes fueron sugeridos por la
retroalimentación docente antes de que sus consecuencias fueran evidentes, y el
tercero emergió de la implementación efectiva del modelo. Esto ilustra que la
coherencia metodológica se verifica al construir, y no solo al declarar.
